\question{9.m1}{Een kritiek gebied $Z \ldots$}
\begin{enumerate}[label=(\alph*)]
    \item bestaat uit alle hypothesen die moeten worden verworpen.
    \item mag geen waarneming bevatten.
    \item geeft voor de toetsingsgrootheid aan welke uitkomsten daarvan tot verwerping van de nulhypothese zullen leiden.
    \item heeft een kans van $1 - \alpha$.
\end{enumerate}
\answer{
    Het kritieke gebied is een manier om een toetsuitslag te bepalen door te kijken naar de kansverdeling van de theoretische toetsingsgrootheid 
    en te bepalen welke uitkomsten daarvan als \enquote{extreem} kunnen worden beschouwd onder de aanname van $H_0$.
    Het juiste antwoord is dus (c).
}
